% Generated by scripts/assemble.py; corrections belong in source/corrections.json.
\examhead{100}{南京航空航天大学 2026 年数学分析试卷(回忆版)}{99}
% p100-b18; source PDF page 100
\blockanchor{p100-b18}
(注.题目来自考生回忆,部分表述或数据可能有误)\par
% p100-b19; source PDF page 100
\blockanchor{p100-b19}
\question{1}计算题.\par
% p100-b20; source PDF page 100
\blockanchor{p100-b20}
(1) 求极限 $\lim_{x\to \infty}\left[\left(x^3 -\frac{1}{2} x^2 -x\right)e^{\frac{1}{x}} - \sqrt{x^6 + 1}\right].$\par
% p100-b21; source PDF page 100
\blockanchor{p100-b21}
(2) 求积分 $\int_{-\frac{\pi}{4}}^{\frac{\pi}{4}} \frac{\tan^2 x}{1 + e^{-x}} \, \mathrm{d}x$.\par
% p100-b22; source PDF page 100
\blockanchor{p100-b22}
\question{2}设 $f(x)$ 在 $[0,1]$ 上可微, $f(1) - 2\int_{0}^{\frac{1}{2}}xf(x)\mathrm{d}x = 0$, 证明: 存在 $\xi \in (0,1)$, 使得 $f'(\xi) = -\frac{f(\xi)}{\xi}$.\par
% p100-b23; source PDF page 100
\blockanchor{p100-b23}
\question{3}压缩数列 $|x_{n} - x_{n - 1}| \leq r|x_{n - 1} - x_{n - 2}|$, $r \in (0,1)$, 证明: $\{x_{n}\}$ 一定收敛.\par
% p101-b00; source PDF page 101
\blockanchor{p101-b00}
\question{4}解答如下问题:\par
% p101-b01; source PDF page 101
\blockanchor{p101-b01}
(1) 未知.\par
% p101-b02; source PDF page 101
\blockanchor{p101-b02}
(2) $f(x) = \frac{x + 2}{x + 1}, a_{1} = 1, a_{n + 1} = f(a_{n})$ ，证明： $\{a_{n}\}$ 收敛，并求极限.\par
% p101-b03; source PDF page 101
\blockanchor{p101-b03}
\question{5}证明导数有界则一致连续, 反之, 一致连续导数是否有界? 证明或给出反例.\par
% p101-b04; source PDF page 101
\blockanchor{p101-b04}
\question{6}解答如下问题:\par
% p101-b05; source PDF page 101
\blockanchor{p101-b05}
(1) 判断  $\int_{1}^{+\infty}\frac{\sin x}{x^{\alpha}+\sin x}dx$  的敛散性.\par
% p101-b06; source PDF page 101
\blockanchor{p101-b06}
(2) 证明: $f(x) = xe^{-x^2} \int_0^x e^{-t^2} \, \mathrm{d}t$ 在 $[0, +\infty)$ 上一致连续.\par
% p101-b07; source PDF page 101
\blockanchor{p101-b07}
\question{7}设 $f(x), g(x)$ 在 $[a, b]$ 上连续.\par
% p101-b08; source PDF page 101
\par\Needspace{6\baselineskip}
\blockanchor{p101-b08}
(1) $[f(x) - f(y)][g(x) - g(y)] \geq 0$, 证明: $\int_{a}^{b} f(x) \, \mathrm{d}x \int_{a}^{b} g(x) \, \mathrm{d}x \geq (b - a) \int_{a}^{b} f(x) g(x) \, \mathrm{d}x$.\par
\nopagebreak[4]
% p101-b09; source PDF page 101
\blockanchor{p101-b09}
\displayformula{(2) \int_ {0} ^ {1} \frac {f (x)}{1 - f (x)} \leq \frac {\int_ {0} ^ {1} f (x) \mathrm{d} x}{\int_ {0} ^ {1} [ 1 - f (x) ] \mathrm{d} x}.}
% p101-b10; source PDF page 101
\blockanchor{p101-b10}
\question{8}证明 $f(x) = \sum_{n=1}^{\infty} \frac{|x - \frac{1}{n}|}{2^n}$ 在 $x_k = \frac{1}{k} (k = 1, 2, \cdots)$ 点不可微，求导函数.\par
% p101-b11; source PDF page 101
\par\Needspace{6\baselineskip}
\blockanchor{p101-b11}
\question{9}已知\par
\nopagebreak[4]
% p101-b12; source PDF page 101
\blockanchor{p101-b12}
\displayformula{f (x, y) = \left\{ \begin{array}{l l} x - y + \frac {x y ^ {2}}{x ^ {2} + y ^ {2}}, & x ^ {2} + y ^ {2} \neq 0, \\ 0, & x ^ {2} + y ^ {2} = 0. \end{array} \right.}
% p101-b13; source PDF page 101
\blockanchor{p101-b13}
求 $f(x,y)$ 偏导存在可微性连续性(???), $f(x,y)$ 在 $(0,0)$ 点沿方向 $\pmb{l} = (\cos \theta, \sin \theta)$ 的方向导数.\par
% p101-b14; source PDF page 101
\blockanchor{p101-b14}
\question{10}$z = z(x,y)$ 存在二阶偏导，$u = x^2 - y^2, v = 2xy, z_{xx} + z_{yy} = 0$，用 $u, v$ 自变量表示.\par
% p101-b15; source PDF page 101
\blockanchor{p101-b15}
\question{11}有一球体半径为 $a$, 球外有一点 $M(0,0,b) (b > a)$, 从点 $M$ 作与球相切的圆锥, 求球与圆锥所围立体的体积.\par
% p101-b16; source PDF page 101
\par\Needspace{6\baselineskip}
\blockanchor{p101-b16}
\question{12}求曲面积分\par
\nopagebreak[4]
% p101-b17; source PDF page 101
\blockanchor{p101-b17}
\displayformula{\iint_ {\Sigma} y ^ {2} z \mathrm{d} x \mathrm{d} y + x z \mathrm{d} x \mathrm{d} z + x y ^ {2} \mathrm{d} y \mathrm{d} z.}
% p101-b18; source PDF page 101
\blockanchor{p101-b18}
$\Sigma$ 为 $x^{2} + y^{2} = z, x^{2} + y^{2} = 1$ 与坐标平面围成立体的表面，取外侧.\par
